% Experiment 3 — why the probe is read at each world's first-response horizon h*.
% GENERATED: paper/figures/exp3_probe_diagnostic.pdf is drawn from the world matrices themselves by
% exp3_training_transfer/dag_family/plot_probe_diagnostic.py, so it cannot drift from the catalog.
% Regenerate with:  python3 plot_probe_diagnostic.py
\begin{center}
\begin{minipage}{\linewidth}
\centering
\includegraphics[width=\linewidth]{figures/exp3_probe_diagnostic.pdf}
\captionof{figure}{\textbf{Why the instrument is read at each world's first-response horizon $h^\star$, and what
counts as success.}
\emph{(a)} The Experiment-2 four-shock probe asks how the poll moves under each of four news
shocks. Read one week ahead it is \emph{degenerate} in any world where news needs longer than a week
to reach the poll: all four scenarios then share a single target, so the answers say nothing about
the hidden gain $g$ and a constant reply is exactly optimal. Six of the twelve worlds are blind this
way. Reading each world at its own $h^\star$ (the earliest week a news pulse moves the
probed poll) restores a $g$-bearing slope in every world where $g$ sits on an input edge ($6/8$ training,
$4/4$ held-out). The two hidden-$\phi$ worlds stay flat by construction: there $g$ is a persistence,
so no input$\to$output response identifies it at any horizon.
\emph{(b)} The consequence for the reward. A constant ``ignore the news'' reply collects $70\%$ of the
oracle's reward under the one-step probe, making a non-inferential policy nearly
optimal. It collects only $46\%$ under $h^\star$.
\emph{(c)} References for the held-out worlds. The oracle is an upper bound. Because the
prompt discloses the structure, the relevant baselines are the \emph{blind constants}: the best
single response $\gamma$ chosen in hindsight per world ($5.57$) and the always-prior $\gamma{=}0.55$
($5.69$). A forecaster that estimates $g$ must land left of them; one that has learned a good
average response cannot.}
\label{fig:exp3-probe}
\end{minipage}
\end{center}
